\documentclass[a5paper,10pt]{article}

% https://github.com/InDevRus/filippov-solutions

% Нумерация страниц
\pagenumbering{gobble}

% Настройка полей
\usepackage{geometry}
\geometry{left=1cm}
\geometry{right=1cm}
\geometry{top=1cm}
\geometry{bottom=1.5cm}

% Вставка таблицы
\usepackage{array}
\usepackage{tabularx}

% Инструменты выравнивания формул
\usepackage{amsmath}
\usepackage{mathtools}

% Диакритика в математике
\usepackage{amsfonts}

% Вычисления размеров на ходу
\usepackage{calc}

% Удобные записи производных
\usepackage[italic=true]{derivative}

% Настройки сносок
\usepackage{enotez}

% Длинные знаки векторов
\usepackage{anyfontsize}
\usepackage[b]{esvect}

% Вставка картинок
\usepackage{graphicx}

% Вставка красной строки
\usepackage{indentfirst}
\setlength{\parindent}{1em}

% Условия в командах
\usepackage{xifthen}

% Луа-скрипты
\usepackage{luacode}

% Настройка команд отдельно в математическом и текстовом режимах
\usepackage{mathcommand}

% Для повторения LaTeX кода
\usepackage{multido}

% Конфигурация отступов формул
\delimitershortfall-1sp
\usepackage{mleftright}
\mleftright

% Растягивание объектов
\usepackage{relsize}

% Обтекание текстом
\usepackage{wrapfig}
\usepackage{subcaption}
\usepackage{floatflt}

% Поддержка ввода кириллицы
\usepackage[english, russian]{babel}

% Настройка корневой позиции для компилятора
\usepackage{import}
\providecommand*{\ProjectRootPath}{..}

\import{\ProjectRootPath/preambles/macros/}{theorems}

\import{\ProjectRootPath/preambles/fonts}{font_setup}

% Знак градуса
\usepackage{siunitx}

\import{\ProjectRootPath/preambles/macros/}{operators}
\import{\ProjectRootPath/preambles/macros/}{redefinitions}

\import{\ProjectRootPath/preambles/fonts}{letters_setup}
\import{\ProjectRootPath/preambles/fonts/adjustments}{custom_superscripts}

\import{\ProjectRootPath/preambles/custom}{formatting}
\import{\ProjectRootPath/preambles/custom}{frames}
\import{\ProjectRootPath/preambles/custom}{list_setup}

% TODO: перевести команды на современный стиль объявления

\begin{document}

Всякая точка экстремума $\sub{x}{0}$ функции $y\br{x}$ такова,
что $y\`\br{\sub{x}{0}} = 0$.
Из этого следует, что уравнение точек экстремума имеет вид $\,f\br{x; y} = 0$.

Условие $\,f\br{x; y} = 0$ является необходимым, но не достаточным, 
так как на кривые, задаваемые уравнением $\,f\br{x; y} = 0$, 
попадут помимо экстремумов в том числе седловые точки.
Отличить такие точки можно следующим образом. Если для 
$\sub{x}{0}$ имеем $f\br{\sub{x}{0}; y} = 0$, 
и при всяком зафиксированном $\sub{y}{0}$ функция 
$g\br{x} = f\br{x; \sub{y}{0}}$ меняет знак при прохождении через 
$\sub{x}{0}$, то имеем точку экстремума, а если не меняет,
то это седловая точка.
Если же имеем экстремум, то при прохождении через 
$\sub{x}{0}$ смена знака с положительного на отрицательный
означает точку максимума, наоборот -- точку минимума.

В случае существования частной производной 
$\parder{f}{x}$ в точке 
$\tsys{x = \sub{x}{0} \\ y = \sub{y}{0}}$
условия выше формулируются так: 
$\parder{f}{x} \br{\sub{x}{0}; \sub{y}{0}} > 0$ соответствует точке минимума;
$\parder{f}{x} \br{\sub{x}{0}; \sub{y}{0}} < 0$ точке максимума. 
Случай $\parder{f}{x} \br{\sub{x}{0}; \sub{y}{0}} = 0$
не даёт ответа на вопрос о наличии и виде экстремума и
требует исследования частных производных высших порядков.

\end{document}
